\section{Conclusion}
\label{sec:conclusion}

MCP tool retrieval is an effective and practical strategy for designing straightforward and
well-defined workflows through the use of an AI agent. RAG performed slightly worse but the
reason for this is not clear and needs further investigation. MCP has a higher token usage due to
the tools list being sent with every request, token caching can alleviate much of this cost if
implemented. This does not reduce the context size, however, and with a significantly larger
number of tools, MCP may degrade while RAG will only retrieve the top N most relevant tools.

Paraphrases did not have a significant impact on the performance of the agentic workflow
generation, showing that the system is robust across a range of similarly worded inputs. Future
work can include a similar test with extreme paraphrases with significantly different semantics
while still containing the same intent and request. Grammatically incorrect prompt inputs can
also enhance the stringency of the tests by checking if the agent can properly correct and
understand the prompt even if the wording is obfuscated. The test can be expanded with additional
scenarios, especially those that are more open-ended and have a less well-defined correct
solution.

A clear prerequisite of the RAG approach is, that the tools are known a priori, as otherwise the
embedding cannot be done to inform the LLM about the available deterministic tools to perform the
tasks. This a priori knowledge is also what saves resources and time during the execution, as the
LLM does not need to query the MCP over and over again ``learning'' about the currently available
models. But if the tooling changes frequently for whatever reason, MCP is the only option.

As more tools of this kind emerge, best practices or tools like benchmarks and test-suites will
be needed to systematically assess an approach's performance on different dimensions. These
dimensions may include the capability to handle paraphrases and other variations of a prompt or
the scalability of the solution in multiple facets such as cost or performance. Ideally, these
assessment tools would allow researchers and developers to answer the question of whether the
implementation of the agentic workflow is appropriate and meaningful for their individual use
cases. Evaluating workflows across different domains will require incorporation of data and API
standards such as OGC processes. We will work towards this in the VelocityAdapt project to better
tackle these questions in the context of urban digital twins.
